(a) The Rees algebra of $\mathcal I^d$ is the $d$th Veronese algebra of the Rees algebra of $\mathcal I$. The isomorphisms of \exref{II.5.13} on affine opens are compatible with restriction, so the blowups are isomorphic over $X$.

(b) Multiplication gives $(\mathcal I\mathcal J)^n\cong\mathcal I^n\otimes\mathcal J^n$, since $\mathcal J$ is an invertible ideal. Thus the two Rees algebras differ by the twist in (7.9), and their Proj schemes are isomorphic.

(c) Choose $\mathcal I$ representing $f$ by (7.17). Since $X$ is regular, the common divisorial factor of $\mathcal I$ is an effective Cartier divisor $D$: its coefficient at a prime divisor is the minimum of the valuations of local generators of $\mathcal I$. Set $\mathcal J=\mathcal I\otimes\mathcal O_X(D)$. Locally in a regular local ring, this divides the generators by their greatest common divisor. Hence $\mathcal J\subseteq\mathcal O_X$ is coherent, has no nonunit common factor, and $\mathcal I=\mathcal J\mathcal O_X(-D)$. Part (b) shows that $\mathcal J$ gives the same blowup.

On $U$, $\mathcal I$ is invertible by (7.13)(a), because its pullback to the blowup is invertible and $f$ is an isomorphism there. Dividing out its common factor gives $\mathcal J|_U=\mathcal O_U$. Conversely, where $\mathcal J=\mathcal O_X$, its blowup is an isomorphism by (7.13)(b), so this open set is contained in $U$. Therefore $V(\mathcal J)=X-U$.
